\documentclass[11pt,a4paper]{article}

% Ein ruhiger Ort für große Gedanken.
\usepackage[utf8]{inputenc}
\usepackage[T1]{fontenc}
\usepackage[ngerman]{babel}
\usepackage{amsmath,amssymb}
\usepackage{graphicx}
\usepackage[margin=2.6cm]{geometry}
\usepackage{hyperref}

\title{Notizen zur Unendlichkeit}
\author{Dein Name}
\date{\today}

\begin{document}
\maketitle

\begin{abstract}
Jede große Idee beginnt mit einem leeren Blatt.
Dies ist dein Raum, um Gedanken zu ordnen,
Zusammenhänge zu entdecken und Neues zu schreiben.
\end{abstract}

\section{Ein Anfang}
Willkommen in \textbf{SimpleLexEditor}. Hier treffen klare
Gedanken auf schönen Satz. Konzentriere dich auf
deine Ideen -- wir kümmern uns um den Rest.

Ein gutes Dokument braucht nicht viel: einen
Gedanken, eine Struktur und die Freiheit,
einfach anzufangen.

\section{Die Schönheit der Mathematik}
Manche Ideen lassen sich am besten in einer
einzigen Zeile ausdrücken. Die Eulersche Identität
verbindet fünf fundamentale Konstanten:

\begin{equation}
  e^{i\pi} + 1 = 0
\end{equation}

Und manchmal eröffnet eine Summe eine ganz
neue Perspektive auf das Unendliche:

\[
  \sum_{n=1}^{\infty} \frac{1}{n^2}
  = \frac{\pi^2}{6}
\]

\section{Was als Nächstes kommt}
\begin{itemize}
  \item Gib deinen Gedanken eine Struktur.
  \item Ergänze Formeln, Tabellen und Abbildungen.
  \item Teile dein fertiges Werk als PDF.
\end{itemize}

\bigskip
\noindent\textit{Der beste Moment anzufangen ist jetzt.}

\end{document}
